\apendice{Plan de proyecto}

\section{Introducción}
Este proyecto aborda el desarrollo de una herramienta web de apoyo a la decisión para la elección de Sistemas Aumentativos y Alternativos de Comunicación (SAAC) en pacientes diagnosticados con Esclerosis Lateral Amiotrófica (ELA). La ELA es una enfermedad neurodegenerativa progresiva que afecta a las neuronas motoras, provocando una pérdida gradual de las capacidades motoras y, con frecuencia, del habla, lo que hace imprescindible el uso de sistemas de comunicación adaptados \cite{lesturnerals:elacomunicacion}.

La selección de un SAAC adecuado es un proceso complejo que depende de múltiples variables clínicas extraídas habitualmente de escalas funcionales como la ALSFRS-R \cite{cedarbaum:alsfrsr}: las capacidades funcionales del paciente, su nivel de alfabetización, el idioma, las plataformas tecnológicas disponibles y el entorno de uso. En la práctica clínica, esta decisión recae habitualmente en logopedas y terapeutas ocupacionales que deben conocer un catálogo amplio y en constante actualización de herramientas.

El objetivo del TFG es desarrollar una aplicación web que, a partir de un cuestionario, aplique un algoritmo de filtrado para proporcionar una recomendación personalizada de sistemas SAAC, periféricos y bancos de voz adaptados al perfil del paciente. Adicionalmente, la herramienta recoge feedback de los usuarios con fines de investigación y cuyo almacenado se realiza de forma anónima en cumplimiento del RGPD.

\section{Metodología de ingeniería del software}
El desarrollo de este proyecto ha seguido una metodología de carácter iterativo e incremental.

Se optó por este enfoque frente a una metodología en cascada por dos razones: Por un lado, los requisitos del sistema evolucionaron a lo largo del proyecto a medida que se fue ampliando el catálogo de sistemas SAAC y se refinó el algoritmo de filtrado. Por otro lado, la naturaleza del proyecto requería validar cada paso antes de poder avanzar al siguiente.

El desarrollo se organizó en las siguientes fases:
\begin{description}
	\item[\textbf{Fase 1: Análisis y requisitos}] Revisión bibliográfica sobre ELA y SAAC, definición de los requisitos funcionales del sistema, y diseño del cuestionario clínico junto con los criterios de filtrado.
	\item[\textbf{Fase 2: Diseño}] Diseño de la arquitectura definida en el anexo C, modelado de la base de datos relacional y definición del algoritmo de filtrado.
	\item[\textbf{Fase 3: Implementación}] Desarrollo iterativo de la aplicación Flask, construcción y rellenado de la base de datos PostgreSQL, implementación del algoritmo de filtrado y desarrollo de las plantillas Jinja2.
	\item[\textbf{Fase 4: Validación}] Pruebas funcionales del cuestionario y del algoritmo, verificación del cumplimiento RGPD y recogida de feedback con casos reales y de prueba.
	\item[\textbf{Fase 5: Documentación}] Redacción de la memoria y los anexos del TFG.
\end{description}

La principal fase que queda fuera del alcance de este TFG y que podría abordarse en un trabajo futuro es la validación clínica formal más extendida que la actual, con una muestra representativa de pacientes y profesionales.

\section{Planificación temporal}
La planificación inicial del proyecto se estableció en febrero de 2026 y se refleja en el siguiente diagrama de Gantt.

\imagen{../img/Planificacion_temporal_inicial.png}{Planificación temporal inicial del proyecto (febrero 2026).}

La planificación seguida finalmente se establece en mayo teniendo en cuenta la evolución de \textit{issues} del GitHub y pequeños cambios de planificación que fueron surgiendo al abordar el proyecto. Se refleja aquí:
 \imagen{../img/gantt_final.png}{Planificación temporal final del proyecto (mayo 2026).}

Dicho repositorio de GitHub se puede consultar a través del siguiente enlace  \url{https://github.com/itslxcy/Trabajo-fin-de-grado.git}

\subsection{Planificación económica}
A continuación se desglosan los costes reales del proyecto, calculados bajo estándares normativos del sector tecnológico y desglosados entre recursos humanos y recursos materiales.

\subsubsection{Recursos humanos}
El proyecto ha sido desarrollado por una única persona en el rol de analista, diseñadora y programadora. De acuerdo con la asignación académica del Trabajo de Fin de Grado, fijada en *2 créditos ECTS, se computa una dedicación oficial de 360 horas de trabajo autónomo.

Para dotar al presupuesto de un rigor de mercado profesional, el cálculo del precio por hora se ha fundamentado de forma estricta en las tablas salariales del XVIII Convenio Colectivo Estatal de Empresas de Consultoría, Tecnologías de la Información y Estudios de Mercado (Área 3: Programación y Gestión de Sistemas, Grupo D: Ingenieros/as Junior).  \cite{boe:conveniotic}

El cálculo detallado del coste por hora se desglosa a continuación:
\begin{itemize}
    \item \textbf{Salario base anual de mercado:} Se toma como referencia un salario bruto de 24.000 €/año para un perfil de ingeniería de software junior en el ámbito de la salud con menos de dos años de experiencia.
    \item \textbf{Coste de cotización patronal:} Aplicando un 31\% en concepto de seguridad social a cargo de la empresa (según lo dispuesto en la Orden PJC/51/2024), el coste real anual para la organización asciende a 31.440 €.
    \item \textbf{Precio hora profesional:} Dividiendo el coste total entre la jornada laboral anual fijada por convenio (1.736 horas), resulta una tarifa horaria de \textbf{18,11 €/h} (redondeada de forma conservadora a 18 €/h).
\end{itemize}

Bajo estas premisas de ingeniería y cumplimiento normativo, el coste de la partida de personal asociado al desarrollo exclusivo de la solución web y su motor lógico es el siguiente:

\begin{table}[!h]
    \centering
    \begin{tabular}{ l p{4cm} r }
        \toprule
        \textbf{Rol} & \textbf{Horas} & \textbf{Coste estimado} \\
        \midrule
        \addlinespace
        Analista / Desarrolladora Junior & 360 h & 6.480 € \\
        \addlinespace
        \bottomrule
    \end{tabular}
    \caption{Coste real de recursos humanos según Convenio TIC.}
\end{table}
\FloatBarrier

Este importe refleja el valor de mercado del esfuerzo de ingeniería necesario para implementar la lógica proposicional del sistema experto, la robustez de las consultas y la capa de seguridad web aplicada bajo estándares OWASP (Flask-Talisman).

\subsubsection{Recursos materiales y software}
Para el desarrollo de la actividad se cuantifica el coste de amortización de los equipos físicos e infraestructura lógica utilizada.

\begin{table}[!h]
    \centering
    % tabularx necesita el ancho total (\textwidth) y usamos 'X' para la columna que se estira automáticamente
    \begin{tabularx}{\textwidth}{ l X r }
        \toprule
        \textbf{Recurso} & \textbf{Descripción} & \textbf{Coste} \\
        \midrule
        \addlinespace
        Ordenador portátil & Equipo de desarrollo & Amortización: 150 € \\
        \addlinespace
        Python / Flask & Framework de desarrollo & Gratuito (\textit{open source}) \\
        \addlinespace
        PostgreSQL / SQLAlchemy & Persistencia y capa ORM & Gratuito (\textit{open source}) \\
        \addlinespace
        Visual Studio Code & Entorno de desarrollo integrado & Gratuito \\
        \addlinespace
        TeXworks & Redacción de la memoria & Gratuito (\textit{open source}) \\
        \addlinespace
        \bottomrule
    \end{tabularx}
    \caption{Coste estimado de recursos materiales y software.}
\end{table}
\FloatBarrier

Considerando ambas partidas, el coste total de inversión del proyecto asciende a 6.630 €, donde el 97,7\% corresponde a la partida de ingeniería de software, justificando la complejidad del motor lógico de recomendación dinámico implementado.

\subsection{Viabilidad legal}
La herramienta desarrollada procesa datos clínicos de carácter sensible, por lo que su diseño se ha ajustado al Reglamento General de Protección de Datos (RGPD, UE 2016/679) \cite{ue:rgpd} y a la Ley Orgánica 3/2018 de Protección de Datos Personales y garantía de los derechos digitales (LOPDGDD) \cite{boe:lopdgdd}.

Las medidas adoptadas para garantizar la conformidad legal son las siguientes:

\begin{itemize}
	\item No se almacena ningún dato de carácter personal que permita la identificación directa del paciente. Los registros de investigación se guardan de forma completamente anónima mediante un identificador autogenerado compuesto por combinaciones aleatorias de letras y números.
	\item El sistema requiere el consentimiento explícito del usuario para el procesamiento obligatorio de datos, sin este no se procesa el cuestionario.
	\item El consentimiento para el almacenamiento con fines de investigación es opcional e independiente del anterior.
	\item La aplicación incluye una página de información legal accesible desde el menú de navegación principal, donde se detalla la finalidad del tratamiento, las bases jurídicas, la política de conservación y los derechos del usuario.
\end{itemize}

\subsection{Análisis de riesgos y limitaciones para la explotación clínica}
Al tratarse de un proyecto de investigación académica con datos completamente anonimizados, el sistema opera en un entorno controlado exento de riesgos de reidentificación. No obstante, de cara a una futura transferencia tecnológica a un entorno sanitario real (como centros de salud o el Sacyl), la viabilidad legal exigiría el cumplimiento del Esquema Nacional de Seguridad \cite{boe:ens}, así como la subsanación de dos requisitos delimitados en \cite{ue:rgpd}:

\begin{enumerate}
    	\item \textbf{Evaluación de Impacto en la Protección de Datos (EIPD):} El tratamiento de datos de salud (categoría especial de datos según el Art. 9 del RGPD) a gran escala obliga a la realización de una EIPD previa según el Art. 35 del RGPD. Este análisis formal de riesgos evaluaría 		el impacto del motor de recomendación sobre los derechos y libertades de los pacientes, identificando medidas de mitigación técnicas antes de cualquier despliegue.
   	 \item \textbf{Designación de un Delegado de Protección de Datos (DPD):} Mientras que en este TFG la supervisión y contacto recae de forma directa en el equipo investigador académico, el despliegue profesional requeriría la tutela del DPD designado legalmente por la entidad 			sanitaria o el hospital explotador, actuando como enlace con la Agencia Española de Protección de Datos (AEPD).
\end{enumerate}